\documentclass{article}
\usepackage{amsmath,hyperref}
\begin{document}
\section*{Distal cave portal evidence beyond the unit-hop model}
\paragraph{Audited domain.}
Take the union of closed one-meter cubes centered at original
lookup coordinates assigned to the 20 selected tiles, as an
explicit center-position authorization envelope. This is not a
claim that the lookup enumerates actual free space, that the
physical cave has no exterior bypass, or that a robot is
authorized to fly there in the stock scenario. Check collision
separately against the pinned triangles of those 20 placed tiles.
The entire development model does not import every world object.

Two closed unit cubes at distinct integral coordinates can touch
only if their coordinate difference has maximum norm at most
one. The upper adjacency graph therefore includes face, edge
and corner contacts regardless of physical collision. The actual
compiled lookup domain contains 138,777 used cubes. Its conservative
tile adjacency has 22 edges, not the 20 of the induced visibility
graph. In particular, there are contacts between 114 and 128 and
between 124 and 157 not represented by the previous travel
abstraction. These contacts are possible bypasses for lower-bound
proof purposes, but not automatically flight paths.

\paragraph{Two full portal cuts and positive motion witnesses.}
All possible cube contacts across $(115,118)$ lie on $x=312.5$,
with $y$ between $-257.5$ and $-242.5$ and $z$ between
$9.5$ and $20.5$. Across $(117,119)$ they lie on $y=-237.5$,
with $x$ between $317.5$ and $333.5$ and $z$ between
$29.5$ and $39.5$. These are bounding rectangles of \emph{all}
possible contacts in the declared union, not measurements of a
particular physical doorway. A specified thin solid slab contains
each corresponding rectangle. Removing either crossing leaves
an alternate in the conservative graph; removing both
disconnects inspection from the start inside this envelope.
Continuous shortcuts touching cube corners cannot bypass that
topological disconnection.

Find actual collision-clear polygonal flights on a one-meter search
grid, then try straight shortcuts using the same imported geometry.
For an integer-endpoint segment, split at every half-integer
plane crossing. A midpoint in every complementary interval
proves the center stays within the closed union, and separate
exact padded swept-box checks prove clearance from every imported
triangle and the named closed slab. Return follows the reverse
segments. Both authored alternatives have full start--inspection--
return witnesses, requiring about 470.63 and 509.02 meters of
total movement. These are constructive upper bounds only, for
the declared body and stay-in-domain law; the calculation is not
a measured execution, a complete trajectory search or a dynamics
certificate.

\paragraph{Status of the earlier information advantage.}
Replacing tile counts by lower bounds on continuous travel
changes the comparison. Enclose the contact set for every
tile-pair in its checked bounding box. Give the robot free choice
of each next crossing point within that box and charge only the
distance between consecutive boxes, rounded down; permit all
transitions within a tile. This is an optimistic metric
relaxation. The two known-world round-trip lower bounds are
348.472 and 246.750 meters, not equal. For local-first
identification from the named endpoint tiles, the same relaxation's
best worst-case lower bound is 348.472 meters. It therefore gives
no certified separation from the known-world upper bounds.
Equality of two lower bounds does not establish equality of
optimal policies. The gap from the unit-hop experiment cannot
be quoted as a result about continuous choice.

Local observations, independently active work and the later
custodian failure required for the research contract are still
separate obligations. This study uses an authored bounded
authorized space, two modeled slabs and partial collision geometry.
It is not an independently selected scenario or I2 acceptance.
\end{document}
